% SPDX-FileCopyrightText: Copyright (c) 2024-2026 Yegor Bugayenko
% SPDX-License-Identifier: MIT

\section{System}\label{sec:system}

All objects presented in this Section belong to the \ff{sm} package.
They provide access to the underlying operating system.

The \deff{os} object represents the current operating system.
The \adeff{name} attribute is its name, while \deff{is-windows},
  \deff{is-linux}, and \deff{is-macos} are predicates about its family.

The \deff{getenv} object reads an environment variable as a \deff{string};
  an empty string means the variable does not exist.

The \deff{eol} object is the
  system-dependent line separator: \ff{\char`\\n} on UNIX
  and \ff{\char`\\r\char`\\n} on Windows.

The \deff{system-clock} object is a wall-clock time source backed by the
  operating system; it decorates an \deff{i64} holding the number of
  milliseconds since the Unix epoch.

The \deff{posix} object makes a Unix system call by name through the POSIX
  interface, while the \deff{win32} object makes a \ff{kernel32.dll}
  function call by name; both accept a \deff{tuple} of arguments.
The call returns an object whose \deff{code} attribute is the numeric
  return code of the syscall and whose \deff{output} attribute is its output.
For example, this code invokes the \ff{getpid} syscall with an empty
  \deff{tuple} of arguments and reads back the identifier
  of the currently running process:

\begin{ffcode}
(sm.posix "getpid" tuple.empty).code
\end{ffcode}
